\section{Week 3: Interface and Interaction Patterns for Human-AI Systems}

\subsection{Topic Summary}

\begin{keybox}[Learning Objectives]
By the end of this week, you will be able to:
\begin{itemize}
  \item \textbf{HAI1}: Identify key human-AI interaction concepts and discuss how these relate to human agency and initiative within the resulting system.
  \item \textbf{HAI2}: Describe the human-AI interaction design process and its distinction from conventional human-centred design.
\end{itemize}
\end{keybox}

\begin{keybox}[Big Picture]
This week moves from theory to practice: how do the cognitive concepts from Week 2 translate into concrete design guidelines?

We cover:
\begin{enumerate}
  \item \textbf{Six usability goals}: Properties that make systems usable.
  \item \textbf{Usability vs.\ User Experience}: Related but distinct concepts.
  \item \textbf{Nielsen's 10 heuristics}: Rules of thumb for achieving usability.
  \item \textbf{AI-specific guidelines}: How heuristics differ for deep learning vs.\ GenAI systems.
  \item \textbf{Building with AI}: Building blocks, workflows, and agents (Anthropic framework).
\end{enumerate}
\end{keybox}

\begin{keybox}[What Examiners Like to Ask]
\begin{itemize}
  \item List and explain the six usability goals; apply them to critique a given system.
  \item Distinguish usability from user experience with examples.
  \item Given a screenshot, identify which Nielsen heuristics are violated and suggest fixes.
  \item Compare 2019 deep-learning guidelines vs.\ 2024 GenAI guidelines.
  \item Explain the difference between building blocks, workflows, and agents.
  \item Connect heuristics back to cognitive concepts (memory, mental models, distributed cognition).
\end{itemize}
\end{keybox}

\subsection{Overview + Core Concepts + Extended Theory}

\subsubsection*{Six Usability Goals}

These goals describe properties of \textbf{usable} systems. Use them to evaluate any design.

\begin{center}
\renewcommand{\arraystretch}{1.4}
\begin{tabular}{@{} l p{8cm} @{}}
\toprule
\textbf{Goal} & \textbf{Key Question} \\
\midrule
1. Effective & Can users accomplish their goals at all? \\
2. Efficient & Can users maintain high productivity once learned? \\
3. Safe & What errors are possible? How easy is recovery? \\
4. Good utility & Does it provide the right functions for user tasks? \\
5. Easy to learn & Can users figure it out by exploring? \\
6. Memorable & Can users remember how to use it after time away? \\
\bottomrule
\end{tabular}
\end{center}

\paragraph*{1. Effective to Use}
\begin{itemize}
  \item \textbf{Question}: Can users learn, work, access information, or buy goods?
  \item A messaging app that can't send messages is not effective.
  \item Even the bad parking sign is \textit{technically} effective---you can eventually figure it out.
\end{itemize}

\paragraph*{2. Efficient to Use}
\begin{itemize}
  \item \textbf{Question}: Once learned, can users sustain high productivity?
  \item A messaging app that sends one letter at a time is effective but not efficient.
  \item The good parking sign is efficient---information jumps out at a glance.
\end{itemize}

\paragraph*{3. Safe to Use}
\begin{itemize}
  \item \textbf{Question}: What errors are possible? How do users recover?
  \item Delete key next to Enter = error-prone design.
  \item ``Are you sure you want to delete?'' dialogs = safety mechanism.
\end{itemize}

\paragraph*{4. Good Utility}
\begin{itemize}
  \item \textbf{Question}: Does it provide appropriate functions for user tasks?
  \item Unix terminal has great utility for developers, zero utility for non-technical users.
  \item A TV remote with 50 buttons has poor utility if users only need 5.
\end{itemize}

\paragraph*{5. Easy to Learn}
\begin{itemize}
  \item \textbf{Question}: Can users work out how to use it by exploring?
  \item Chat interfaces are easy to \textit{start} (type and send), but hard to \textit{master} (what can it do?).
  \item Games with clear mechanics are often easier to learn than multipurpose tools.
\end{itemize}

\paragraph*{6. Memorable (Easy to Remember)}
\begin{itemize}
  \item \textbf{Question}: What support helps users remember infrequent tasks?
  \item Links to Week 2: recognition easier than recall; distributed cognition (leave traces).
  \item Tooltips, recent actions, undo history all support memorability.
\end{itemize}

\subsubsection*{Usability vs.\ User Experience}

\begin{defbox}[User Experience (UX)]
How the system \textbf{feels} to the user---their subjective experience of interacting with it.

\textbf{Key insight}: You can't design \textit{a} user experience. You can only design \textit{for} user experience.
\end{defbox}

\paragraph*{Comparison}

\begin{center}
\renewcommand{\arraystretch}{1.3}
\begin{tabular}{@{} l p{5cm} p{5cm} @{}}
\toprule
& \textbf{Usability} & \textbf{User Experience} \\
\midrule
Focus & Task completion & How it feels \\
Scope & Narrow (effective, efficient, safe) & Broad (enjoyable, motivating, beautiful) \\
Measurable & More objective & More subjective \\
\bottomrule
\end{tabular}
\end{center}

\paragraph*{``Beautiful Things Work Better''---But Do They?}

\begin{exbox}[The Juicy Salif Paradox]
\textbf{Philippe Starck's Juicy Salif} (£200): Beautiful sculpture, terrible juicer (sprays pips, makes mess).

\textbf{Argos Value Juicer} (99p): Ugly, but actually juices effectively.

The Juicy Salif has great aesthetics but poor usability. The enjoyment comes from \textit{looking at it}, not \textit{using it}.
\end{exbox}

\paragraph*{The 2×2 Matrix}

\begin{center}
\renewcommand{\arraystretch}{1.3}
\begin{tabular}{@{} l | c c @{}}
& \textbf{High Usability} & \textbf{Low Usability} \\
\midrule
\textbf{High UX} & Beautiful \& usable (goal!) & Beautiful but unusable (Juicy Salif) \\
\textbf{Low UX} & ``Charming'' but usable & What nobody wants \\
\end{tabular}
\end{center}

\subsubsection*{What is a Heuristic?}

\begin{defbox}[Heuristic]
A strategy derived from previous experience with similar problems---a \textbf{rule of thumb} that often works but isn't guaranteed.

In psychology: Simple, efficient rules (learned or evolved) for making decisions under uncertainty or incomplete information.
\end{defbox}

\paragraph*{Key Concepts}
\begin{itemize}
  \item \textbf{Satisficing} (Herbert Simon): Users seek solutions that are ``good enough,'' not optimal.
  \item \textbf{Bounded rationality}: We can't process all information, so we use shortcuts.
  \item Users don't carefully analyze your system---they make quick judgments.
\end{itemize}

\paragraph*{Two Types of Heuristics in Design}
\begin{enumerate}
  \item \textbf{Design heuristics}: How designers solve problems (e.g., ``use consistent icons'').
  \item \textbf{User heuristics}: How users make decisions (e.g., ``this looks clickable'').
\end{enumerate}

\subsubsection*{Nielsen's 10 Usability Heuristics}

Jakob Nielsen's heuristics are \textbf{design principles}---if followed, they increase the likelihood of a usable system.

\begin{center}
\renewcommand{\arraystretch}{1.3}
\begin{tabular}{@{} c l p{7cm} @{}}
\toprule
\textbf{\#} & \textbf{Heuristic} & \textbf{Summary} \\
\midrule
1 & Visibility \& Feedback & Keep users informed about system status. \\
2 & Match Real World & Use user language, not system jargon. \\
3 & User Control \& Freedom & Provide undo, redo, emergency exits. \\
4 & Consistency \& Standards & Same action = same representation. \\
5 & Error Prevention & Prevent errors before they happen. \\
6 & Recognition > Recall & Make options visible; don't rely on memory. \\
7 & Flexibility \& Efficiency & Support both novices and experts. \\
8 & Aesthetic \& Minimalist & High signal-to-noise ratio. \\
9 & Error Recovery & Clear error messages; constructive solutions. \\
10 & Help \& Documentation & Easy to search, task-focused, concrete steps. \\
\bottomrule
\end{tabular}
\end{center}

\paragraph*{Heuristic 1: Visibility \& Feedback}
\begin{itemize}
  \item System should always keep users informed about what's happening.
  \item \textbf{Example}: Selected items change color; Mac charge cable shows orange/green.
  \item \textbf{GenAI}: Perplexity shows ``thinking'' steps---visibility of internal process.
\end{itemize}

\paragraph*{Heuristic 2: Match Between System \& Real World}
\begin{itemize}
  \item Use words, phrases, concepts familiar to users---not system-oriented terms.
  \item \textbf{Skeuomorphic design}: Digital highlighter looks like physical highlighter.
  \item \textbf{Caution}: Chat interfaces trigger human conversation mental models---can lead to wrong expectations about AI capabilities.
\end{itemize}

\paragraph*{Heuristic 3: User Control \& Freedom}
\begin{itemize}
  \item Users make mistakes---provide ``emergency exits'' (undo, cancel, go back).
  \item \textbf{Bad}: Irreversible action after user ignores a dialog.
  \item \textbf{GenAI}: Claude allows branching conversations, editing prompts, regenerating responses.
\end{itemize}

\paragraph*{Heuristic 4: Consistency \& Standards}
\begin{itemize}
  \item Same action = same representation, both within your app and across the ecosystem.
  \item Apple's Human Interface Guidelines: Specific menu orders, icon styles.
  \item \textbf{AI guidelines}: More high-level (``support efficient dismissal'') than prescriptive.
\end{itemize}

\paragraph*{Heuristic 5: Error Prevention}
\begin{itemize}
  \item Better than good error messages is preventing errors entirely.
  \item \textbf{Poka-yoke} (Japanese): Mistake-proofing by design.
  \item \textbf{Example}: Manual car won't start unless clutch is depressed.
  \item \textbf{UI}: Dangerous action buttons are secondary; confirmation dialogs; disabled states.
\end{itemize}

\begin{tipbox}[Slips vs.\ Mistakes]
\textbf{Slip}: User intends correct action, executes wrong one (pressing Delete instead of Enter).

\textbf{Mistake}: User forms wrong intention (wrong mental model of thermostat).

Error prevention addresses both.
\end{tipbox}

\paragraph*{Heuristic 6: Recognition Rather Than Recall}
\begin{itemize}
  \item Minimize memory load: make options, actions, objects visible.
  \item \textbf{Links to Week 2}: Recognition is easier than recall.
  \item \textbf{Examples}: Autocomplete, recently used items, visible toolbars.
\end{itemize}

\paragraph*{Heuristic 7: Flexibility \& Efficiency of Use}
\begin{itemize}
  \item Accelerators for experts; guided paths for novices.
  \item \textbf{Perpetual intermediates}: Most users are neither beginners nor experts.
  \item \textbf{Examples}: Keyboard shortcuts, customizable toolbars, Adobe shortcut stickers.
\end{itemize}

\paragraph*{Heuristic 8: Aesthetic \& Minimalist Design}
\begin{itemize}
  \item Every extra element competes with relevant information.
  \item \textbf{Signal-to-noise ratio}: Maximize signal, minimize noise.
  \item \textbf{Example}: 1990s cluttered websites vs.\ modern landing pages with clear CTAs.
\end{itemize}

\paragraph*{Heuristic 9: Help Users Recognize, Diagnose, Recover from Errors}
\begin{itemize}
  \item Error messages in plain language, precisely indicate problem, suggest solution.
  \item \textbf{Example}: Twitter character counter shows exactly how many characters over.
  \item \textbf{GenAI}: Retry buttons, model switching, conversation branching.
\end{itemize}

\paragraph*{Heuristic 10: Help \& Documentation}
\begin{itemize}
  \item Even if system is usable without docs, provide searchable, task-focused help.
  \item \textbf{Good help}: Concrete steps, not too large, focused on user tasks.
\end{itemize}

\subsubsection*{Heuristics for AI Systems}

The field is developing AI-specific guidelines. Key comparison:

\paragraph*{2019 Microsoft Guidelines (Deep Learning)}
\begin{itemize}
  \item Focused on making AI behavior \textbf{predictable and consistent}.
  \item Assumption: AI systems should behave deterministically.
  \item Example guideline: ``Make clear why the system did what it did'' (G11).
  \item Mental model: AI produces single output that is right or wrong.
\end{itemize}

\paragraph*{2024 IBM Guidelines (Generative AI)}
\begin{itemize}
  \item Embraces \textbf{generative variability} as a fundamental characteristic.
  \item Helps users \textbf{manage and leverage} the fact that GenAI produces different outputs for same input.
  \item Example guideline: ``Leverage multiple outputs,'' ``Help user craft effective outcome specifications.''
  \item Mental model: AI enables \textbf{exploration of possibility spaces}.
\end{itemize}

\paragraph*{Key Differences}

\begin{center}
\renewcommand{\arraystretch}{1.3}
\begin{tabular}{@{} l p{5cm} p{5cm} @{}}
\toprule
& \textbf{2019 Deep Learning} & \textbf{2024 Generative AI} \\
\midrule
Outputs & Single recommendation (accept/reject) & Multiple possibilities (explore) \\
Variability & Try to eliminate & Embrace and leverage \\
User role & Judge correctness & Craft specifications, explore \\
Key guideline & ``Support efficient correction'' & ``Help craft outcome specifications'' \\
\bottomrule
\end{tabular}
\end{center}

\subsubsection*{Building with AI: Anthropic Framework}

From Anthropic's ``Building Effective Agents'' blog---required reading.

\paragraph*{Three Levels of Complexity}

\begin{enumerate}
  \item \textbf{Building Blocks}: Single LLM call with optional tool use.
  \begin{itemize}
    \item Input $\to$ LLM $\to$ Output
    \item Simplest unit; deterministic path.
  \end{itemize}
  
  \item \textbf{Workflows}: Predetermined paths with multiple LLM calls.
  \begin{itemize}
    \item Sequence of blocks with conditional logic.
    \item No cycles; linear or branching.
    \item Example: Input $\to$ LLM1 $\to$ Check $\to$ LLM2 $\to$ Output
  \end{itemize}
  
  \item \textbf{Agents}: LLM controls its own reasoning and tool use.
  \begin{itemize}
    \item Has feedback loops---LLM decides what to do next.
    \item Reasoning and planning; not fully predetermined.
    \item Example: LLM $\to$ Plan $\to$ Act $\to$ Observe $\to$ (repeat)
  \end{itemize}
\end{enumerate}

\begin{tipbox}[Key Distinction]
\textbf{Workflows}: Human designs the control flow; LLM executes steps.

\textbf{Agents}: LLM decides the control flow; human sets goals/constraints.
\end{tipbox}

\subsubsection*{Connecting Heuristics to Usability Goals}

The 10 heuristics help achieve the 6 usability goals:

\begin{center}
\renewcommand{\arraystretch}{1.2}
\begin{tabular}{@{} l l @{}}
\toprule
\textbf{Usability Goal} & \textbf{Related Heuristics} \\
\midrule
Effective & Match real world, Good utility \\
Efficient & Flexibility \& efficiency, Aesthetic \& minimal \\
Safe & Error prevention, Error recovery, User control \\
Good utility & Match real world, Help \& documentation \\
Easy to learn & Recognition > recall, Consistency, Visibility \\
Memorable & Recognition > recall, Consistency \\
\bottomrule
\end{tabular}
\end{center}

\subsubsection*{Worked Examples}

\begin{exbox}[Heuristic Evaluation]
\textbf{Scenario}: A mobile banking app shows ``Error 0x4F2A'' when login fails.

\textbf{Violated heuristics}:
\begin{itemize}
  \item \#2 Match real world: System jargon, not user language.
  \item \#9 Error recovery: No constructive suggestion for fixing the problem.
\end{itemize}

\textbf{Fix}: ``Incorrect password. Please try again or reset your password.''
\end{exbox}

\begin{exbox}[Comparing AI Guidelines]
\textbf{Scenario}: You're designing a writing assistant.

\textbf{2019 approach} (deep learning): Show single ``best'' suggestion; user accepts or rejects.

\textbf{2024 approach} (GenAI): Show multiple variations; let user explore, combine, iterate; help user specify what ``good'' means for their context.
\end{exbox}

\begin{exbox}[Building Block vs.\ Workflow vs.\ Agent]
\textbf{Building block}: ``Summarize this document'' $\to$ LLM $\to$ Summary.

\textbf{Workflow}: Document $\to$ Extract key points (LLM1) $\to$ Check for completeness (rule) $\to$ Generate summary (LLM2) $\to$ Summary.

\textbf{Agent}: ``Research this topic and write a report'' $\to$ Agent plans steps, searches web, reads sources, decides when to stop, writes report.
\end{exbox}

\subsubsection*{Key Definitions Summary}

\begin{center}
\renewcommand{\arraystretch}{1.3}
\begin{tabular}{@{} l p{8cm} @{}}
\toprule
\textbf{Term} & \textbf{Definition} \\
\midrule
Usability & How effectively, efficiently, and satisfyingly users achieve goals \\
User Experience & How the system feels to the user (subjective) \\
Heuristic & Rule of thumb derived from experience; not guaranteed \\
Satisficing & Seeking ``good enough'' rather than optimal solutions \\
Poka-yoke & Mistake-proofing by design \\
Signal-to-noise & Ratio of useful information to clutter \\
Building block & Single LLM call (simplest unit) \\
Workflow & Predetermined sequence of LLM calls \\
Agent & LLM with feedback loop controlling its own actions \\
\bottomrule
\end{tabular}
\end{center}

\subsubsection*{Recommended Reading}
\begin{itemize}
  \item \textbf{Textbook}: Section 1.7 (Usability goals)
  \item \textbf{Blog}: Anthropic---``Building Effective Agents'' (building blocks, workflows, agents)
  \item \textbf{Article}: ``Human-Centered AI: 5 Key Frameworks for UX Designers''
  \item \textbf{Pattern library}: Shape of AI (\texttt{shapeof.ai})
  \item \textbf{Critical thinking}: MIT Critical Thinking paper, Bloom's Taxonomy
\end{itemize}


